\documentclass[10pt,a4paper]{amsart}
\usepackage[margin=19mm]{geometry}
\usepackage{amsmath,amssymb,mathtools}
\usepackage[scr=rsfs,cal=euler]{mathalfa}
\usepackage{xfrac}
\usepackage[unicode,hidelinks,bookmarks=false]{hyperref}
\newcommand{\ie}{i.e.\ }
\newcommand{\cf}{cf.\ }
\newcommand{\resp}{respectively\ }
\newcommand{\Subsection}{\subsection}
\providecommand{\nobreakdash}{\nobreak\hbox{-}}
\setlength{\emergencystretch}{2em}
\multlinegap0pt
\usepackage{tikz}
\usetikzlibrary{arrows.meta,calc,positioning,shapes.geometric}
\allowdisplaybreaks
\newtheorem{theorem}{Theorem}[section]
\newtheorem{lemma}[theorem]{Lemma}
\newtheorem{corollary}[theorem]{Corollary}
\newtheorem{problem}[theorem]{Problem}
\newtheorem{conjecture}[theorem]{Conjecture}
\newcommand{\chik}{\chi_k'}
\begin{document}
\title[Counterexamples to two conjectures on modular edge colorings]{Counterexamples to two conjectures on modular edge colorings of graphs}
\author{Urumqi, Xinjiang 830046, P. R. China; Chunqiang Guo, Baoyindureng Wu}
\date{}
\maketitle
\noindent\emph{Source and licence.} Counterexamples to two conjectures on modular edge colorings of graphs. Version arXiv:2608.10687v1 (CC BY 4.0). This material is distributed under the Creative Commons Attribution 4.0 International licence, http://creativecommons.org/licenses/by/4.0/, as recorded in the arXiv OAI licence metadata for this exact version and shown on the abstract page https://arxiv.org/abs/2608.10687v1. Full rights notice: licenses/arxiv-2608.10687-CC-BY-4.0.txt.\par\smallskip
\noindent\textit{Editorial scope.} Counted unit: the original \texttt{theorem} environment \texttt{thm1} at source lines 154--165 together with its complete proof at lines 327--337; the delivered document keeps source lines 75--338 so that every referenced label and every citation resolves inside it. Native editable source is the paper's own concatenated TeX; the AMS wrapper is editorial formatting only. No publisher class, upstream script or unrelated artwork is redistributed.
\section{Introduction}\label{sec:introduction}

Throughout the paper, $k\geq2$ is an integer, and all graphs are finite
and simple. A \emph{$\chi_k'$-coloring} of a graph $G$ is an edge-coloring in
which the subgraph spanned by the edges of each color has all nonzero
degrees congruent to $1\pmod{k}$. The \emph{mod $k$ chromatic index}
$\chi_k'(G)$ is the minimum number of colors in such a coloring. A graph
$G$ is a \emph{$0_k$-graph} if $d_G(v)\equiv0\pmod{k}$ for every
$v\in V(G)$.

The case $k=2$ is the odd edge-coloring problem. In this setting, the
edges of each color span an odd subgraph, meaning that every nonisolated
vertex of the corresponding color subgraph has odd degree.
Pyber~\cite{Pyber1992} initiated the study of the odd chromatic index and
proved that every simple graph has odd chromatic index at most four. Odd
edge-colorings have since been studied for several graph classes and for
loopless multigraphs; see, for example,
\cite{LuzarPetrusevskiSkrekovski2015,AtanasovPetrusevskiSkrekovski2016,
Petrusevski2018}. Related decomposition problems for odd subgraphs and
matchings were investigated by Kano and Katona~\cite{KanoKatona2002},
and M\'atrai~\cite{Matrai2006} studied coverings by three odd subgraphs.
Thus the definition above extends the parity condition from modulus two
to arbitrary modulus.

Pyber~\cite{Pyber1992} asked whether $\chik(G)$ is bounded by a function
of $k$ alone. Scott~\cite{Scott1997} answered this question affirmatively
by showing that $\chik(G)\leq5k^2\log k$ for every graph $G$, and asked
whether a linear bound in $k$ always holds. Botler, Colucci, and
Kohayakawa~\cite{BotlerColucciKohayakawa2023} proved
$\chik(G)\leq198k-101$ and proposed the following stronger additive
form, stated as Conjecture~1 in their paper.

\begin{conjecture}[Botler, Colucci, and Kohayakawa~\cite{BotlerColucciKohayakawa2023}]
\label{conj:additive-constant}
There exists an absolute constant $C$ such that
$\chik(G)\leq k+C$ for every integer $k\geq2$ and every graph $G$.
\end{conjecture}

Their work on random graphs also showed that the mod $k$ chromatic index
is close to its natural lower bound $k$ with high probability over a
wide range of edge probabilities~\cite{BotlerColucciKohayakawaRandom2023}.
Nweit and Yang~\cite{NweitYang2024} subsequently improved the general
linear bound to $177k-93$. Berthe et al.~\cite{BertheEtAl2026} later proved
$\chik(G)\leq9k+o(k)$ for all $k$ and $\chik(G)\leq7k+o(k)$ when $k$ is
odd. Their method, as in the earlier linear-bound arguments, is closely
related to divisible subgraphs and graph factors with prescribed
degrees modulo $k$; see also
\cite{AlonFriedlandKalai1984,Mader1972,Thomassen2014}. In their concluding
remarks, they proposed the following conjecture, which is Conjecture~13
in their paper.

\begin{conjecture}[Berthe et al.~\cite{BertheEtAl2026}]\label{conj:berthe}
There exists a function $f(k)=o(k)$ such that every $0_k$-graph $G$
satisfies $\chik(G)\leq k+f(k)$.
\end{conjecture}

The conjecture is natural because, in any $\chi_k'$-coloring of a
$0_k$-graph, each color appearing at a vertex contributes $1$ modulo
$k$ to its degree. Hence, at every nonisolated vertex, the number of
colors that appear is a positive multiple of $k$ and is therefore at
least $k$. The conjecture asks whether the whole graph can always be
colored using only slightly more than this local lower bound.

Very recently, Liu, Xu, and Yang~\cite{LiuXuYang2026} disproved
Conjecture~\ref{conj:additive-constant}, even for bipartite graphs.
For all sufficiently large $k$, they constructed bipartite graphs $G$
with
\[
\chi_k'(G)\geq \frac32 k-O((k\log k)^{1/3}).
\]
Their constructions are not $0_k$-graphs, and the codegree obstruction
used in their proofs does not directly apply under the requirement that
every vertex degree be divisible by $k$. Thus their result does not
settle Conjecture~\ref{conj:berthe}.

We show that Conjecture~\ref{conj:berthe} is false even under the
stronger restriction that every vertex degree belongs to $\{k,2k\}$.
Our main result is the following.


\begin{theorem}\label{thm1}
There is a sequence of integers $t_k$, with $1\leq t_k\leq k-1$,
such that the following holds. Let $\{G_k\}_{k\geq2}$ be a sequence
of graphs, where $G_k$ has a vertex partition
$V(G_k)=X_k\cup Y_k$ with $|X_k|=2k-t_k$ and $|Y_k|=2k$.
Suppose that $A_k\subseteq X_k$ has size $t_k$, every vertex of
$A_k$ has degree $2k$, every vertex of $V(G_k)\setminus A_k$ has
degree $k$, and $e(G_k[X_k])=o(k^2)$. Then
\[
\chik(G_k)\geq\left(4-2\sqrt2+o(1)\right)k.
\]
\end{theorem}


\begin{corollary}\label{cor:counterexamples}
Both Conjectures~\ref{conj:additive-constant} and~\ref{conj:berthe} are
false even for connected bipartite $0_k$-graphs with degree set
$\{k,2k\}$. They are also false for connected nonbipartite $0_k$-graphs
with the same degree set.
\end{corollary}

\section{Preliminaries}\label{sec:notation}

We use standard graph-theoretic notation. For a graph $G$, let $d_G(v)$
denote the degree of $v$, and let $G[S]$ be the subgraph induced by
$S\subseteq V(G)$. For two disjoint vertex sets $S$ and $T$, we say
that $S$ is \emph{complete to} $T$ if every vertex of $S$ is adjacent
to every vertex of $T$. A vertex of degree $j$ is called a $j$-vertex,
and $[q]=\{1,2,\ldots,q\}$.

Let $\varphi:E(G)\to[q]$ be a $\chi_k'$-coloring of $G$. For
$c\in[q]$, let $E_c=\{e\in E(G):\varphi(e)=c\}$ and
$G_c=(V(G),E_c)$. Thus $G_c$ is the spanning subgraph formed by the
edges of color $c$. We call $d_{G_c}(v)$ the \emph{color degree} of
$c$ at $v$, and say that $c$ \emph{appears at} $v$ if
$d_{G_c}(v)>0$. Every positive color degree is congruent to
$1\pmod{k}$.

\begin{lemma}\label{lem:k-colors-at-vertex}
Let $G$ be a $0_k$-graph with no isolated vertices, and suppose that
$G$ has a $\chi_k'$-coloring with fewer than $2k$ colors. Then exactly
$k$ colors appear at every vertex. Moreover, at a $k$-vertex $v$, every
color appearing at $v$ has color degree $1$. At a $2k$-vertex $v$, one
color appearing at $v$ has color degree $k+1$, and each of the other
$k-1$ colors appearing at $v$ has color degree $1$.
\end{lemma}

\begin{proof}
For a vertex $v$, let $r(v)$ be the number of colors that appear at
$v$. Each positive color degree is congruent to $1$ modulo $k$, while
$d_G(v)\equiv0\pmod{k}$. Hence $r(v)\equiv0\pmod{k}$. Since $v$ is not
isolated and fewer than $2k$ colors are used, $1\leq r(v)<2k$, so
$r(v)=k$.

If $d_G(v)=k$, the $k$ positive color degrees sum to $k$, so all are
$1$. If $d_G(v)=2k$, each positive color degree is congruent to $1$
modulo $k$ and is at most $2k$, so it is either $1$ or $k+1$. Starting
with $k$ contributions of $1$ leaves a total of $k$, so exactly one
color has color degree $k+1$ and all the others have color degree $1$.
\end{proof}

Whenever a $\chi_k'$-coloring with fewer than $2k$ colors is fixed, the
unique color having color degree $k+1$ at a $2k$-vertex is called the
\emph{heavy color} at that vertex.

\section{A lower bound}\label{sec:lower-bound}

Let $1\leq t\leq k-1$. We consider graphs with a vertex partition
$X\cup Y$ of sizes $2k-t$ and $2k$. The sets $X$ and $Y$ need not be
independent.

\begin{lemma}\label{lem:general-lower-bound}
Assume that $G$ has a vertex partition $V(G)=X\cup Y$ with $|X|=2k-t$ and
$|Y|=2k$. Suppose that $A\subseteq X$ has size $t$, every vertex of
$A$ has degree $2k$, and every vertex of $V(G)\setminus A$ has degree
$k$. Set $p_X=e(G[X])$. Then
\begin{equation}\label{eq:general-bound}
\chik(G)\geq
k+\max\left\{0,
\left\lceil\frac{t(k-t)-2p_X}{2k-t}\right\rceil
\right\}.
\end{equation}
\end{lemma}

\begin{proof}
Since $p_X\geq0$ and $t(k-t)/(2k-t)<t$, the term added to $k$ in
\eqref{eq:general-bound} is at most $t\leq k-1$. Thus there is nothing
to prove if $\chik(G)\geq2k$.

Suppose that $\chik(G)<2k$. Let $q=\chik(G)$, fix a
$\chi_k'$-coloring with color set $[q]$, and let $q=k+s$. The degree
assumptions make $G$ a $0_k$-graph with no isolated vertices, so
Lemma~\ref{lem:k-colors-at-vertex} shows that exactly $k$ colors appear
at every vertex and $s\geq0$.

For $c\in[q]$, let $N_c^X$ and $N_c^Y$ be the numbers of vertices in
$X$ and $Y$, respectively, at which $c$ appears. Let $\mu_c$ be the
number of vertices in $A$ whose heavy color is $c$, and put
$e_c^X=e(G_c[X])$ and $e_c^Y=e(G_c[Y])$. Comparing the two degree sums
for color $c$ gives
\[
N_c^X+k\mu_c-N_c^Y=2(e_c^X-e_c^Y).
\]
Indeed, each appearance in $X\setminus A$ or $Y$ contributes $1$, while
a heavy color at a vertex of $A$ contributes an additional $k$. Edges
between $X$ and $Y$ contribute once to each degree sum and cancel,
whereas edges inside $X$ and $Y$ contribute twice.

Let $\mathcal C_h=\{c\in[q]:\mu_c>0\}$ and
$q_h=|\mathcal C_h|$. Since every vertex of $A$ has exactly one heavy
color, $\sum_{c\in\mathcal C_h}\mu_c=t$ and $q_h\leq t$. If
$c\notin\mathcal C_h$, then $N_c^X\leq2k-t$. If
$c\in\mathcal C_h$, the preceding identity, together with
$N_c^Y\leq2k$ and $e_c^Y\geq0$, gives
$N_c^X\leq2k-k\mu_c+2e_c^X$.

Counting the pairs $(v,c)$ with $v\in X$ at which $c$ appears gives
\begin{align*}
k(2k-t)=\sum_{c=1}^{q}N_c^X
&\leq (q-q_h)(2k-t)
 +\sum_{c\in\mathcal C_h}(2k-k\mu_c+2e_c^X)\\
&=q(2k-t)+q_ht-kt
 +2\sum_{c\in\mathcal C_h}e_c^X\\
&\leq q(2k-t)-t(k-t)+2p_X.
\end{align*}
The last inequality uses $q_h\leq t$ and
$\sum_{c\in\mathcal C_h}e_c^X\leq p_X$. Substituting $q=k+s$ gives
$s(2k-t)\geq t(k-t)-2p_X$. Since $s$ is a nonnegative integer,
\eqref{eq:general-bound} follows.
\end{proof}

The degree sums over $X$ and $Y$ are both $2k^2$. If $m$ is the number
of edges between $X$ and $Y$, then
$2k^2=m+2e(G[X])=m+2e(G[Y])$. Hence
$e(G[X])=e(G[Y])$. Thus the same number of internal edges occurs on the
two sides of the partition.

\subsection{Choosing \texorpdfstring{$t$}{t}}
Let $g_k(t)=t(k-t)/(2k-t)$ for $1\leq t\leq k-1$.

\begin{lemma}\label{lem:choice-t}
For $k\geq2$, let
$\theta_k=(4k-1-\sqrt{8k^2+1})/2$ and
$t_k=\lceil\theta_k\rceil$. Then $1\leq t_k\leq k-1$ and
$g_k(t_k)=\max_{1\leq t\leq k-1}g_k(t)$. Moreover,
\[
\frac{t_k}{k}\longrightarrow2-\sqrt2
\qquad\text{and}\qquad
\frac{g_k(t_k)}{k}\longrightarrow3-2\sqrt2.
\]
\end{lemma}

\begin{proof}
For $k=2$, the only admissible value of $t$ is $1$, and
$0<\theta_2<1$, so $t_2=1$ and the assertion is immediate.
Hence assume $k\geq3$. For
$1\leq t\leq k-2$, direct calculation gives
\[
g_k(t+1)-g_k(t)=
\frac{t^2+(1-4k)t+2k^2-2k}{(2k-t)(2k-t-1)}.
\]
The denominator is positive. The roots of the numerator are $\theta_k$
and $(4k-1+\sqrt{8k^2+1})/2$, and the second root is larger than
$k-1$. Hence the difference is positive for integers $t<\theta_k$ and
negative for integers $t>\theta_k$; if $\theta_k$ is an integer, it is
zero there. This proves the assertion about $g_k(t_k)$.

Since $(4k-1)^2-(8k^2+1)=8k(k-1)>0$ and
$8k^2+1-(2k+1)^2=4k(k-1)>0$, we have $0<\theta_k<k-1$. Finally,
$\theta_k/k\to2-\sqrt2$, so $t_k/k\to2-\sqrt2$. Since
$g_k(t)/k=(t/k)(1-t/k)/(2-t/k)$, we obtain
$g_k(t_k)/k\to3-2\sqrt2$.
\end{proof}


\begin{proof}[Proof of Theorem~\ref{thm1}]
Apply Lemma~\ref{lem:general-lower-bound} with $t=t_k$ and
$p_X=e(G_k[X_k])$. By Lemma~\ref{lem:choice-t}, we have 
$g_k(t_k)=(3-2\sqrt2+o(1))k$. Also,
$2p_X/(2k-t_k)=o(k)$ because $p_X=o(k^2)$ and
$t_k/k\to2-\sqrt2$. Thus the quantity inside the ceiling in
\eqref{eq:general-bound} is positive for all sufficiently large $k$, and
\[
\chik(G_k)\geq(4-2\sqrt2+o(1))k.
\]
\end{proof}


\begin{thebibliography}{99}
\bibitem[AlonFriedlandKalai1984]{AlonFriedlandKalai1984} N. Alon, S. Friedland, G. Kalai, Regular subgraphs of almost regular graphs, J. Combin. Theory Ser. B 37 (1984) 79--91.
\bibitem[AtanasovPetrusevskiSkrekovski2016]{AtanasovPetrusevskiSkrekovski2016} R. Atanasov, M. Petru\v{s}evski, R. \v{S}krekovski, Odd edge-colorability of subcubic graphs, Ars Math. Contemp. 10 (2016) 359--370.
\bibitem[BertheEtAl2026]{BertheEtAl2026} G. Berthe, M. Bonamy, F. Botler, G. Carenini, L. Colucci, A. Dumas, F. Ghasemi, P. M. Viana Neto, On modular edge colorings of graphs, SIAM J. Discrete Math. 40 (2026) 897--904.
\bibitem[BotlerColucciKohayakawa2023]{BotlerColucciKohayakawa2023} F. Botler, L. Colucci, Y. Kohayakawa, The mod $k$ chromatic index of graphs is $O(k)$, J. Graph Theory 102 (2023) 197--200.
\bibitem[BotlerColucciKohayakawaRandom2023]{BotlerColucciKohayakawaRandom2023} F. Botler, L. Colucci, Y. Kohayakawa, The mod $k$ chromatic index of random graphs, J. Graph Theory 103 (2023) 767--779.
\bibitem[KanoKatona2002]{KanoKatona2002} M. Kano, G. Y. Katona, Odd subgraphs and matchings, Discrete Math. 250 (2002) 265--272.
\bibitem[LiuXuYang2026]{LiuXuYang2026} X.-C. Liu, B. Xu, X. Yang, Linear lower bounds for the modular chromatic index, arXiv:2608.02239 (2026).
\bibitem[LuzarPetrusevskiSkrekovski2015]{LuzarPetrusevskiSkrekovski2015} B. Lu\v{z}ar, M. Petru\v{s}evski, R. \v{S}krekovski, Odd edge coloring of graphs, Ars Math. Contemp. 9 (2015) 277--287.
\bibitem[Mader1972]{Mader1972} W. Mader, Existenz $n$-fach zusammenh\"an\-gender Teil\-graphen in Graphen gen\"u\-gend grosser Kanten\-dichte, Abh. Math. Semin. Univ. Hambg. 37 (1972) 86--97.
\bibitem[Matrai2006]{Matrai2006} T. M\'atrai, Covering the edges of a graph by three odd subgraphs, J. Graph Theory 53 (2006) 75--82.
\bibitem[NweitYang2024]{NweitYang2024} O. Nweit, D. Yang, On the mod $k$ chromatic index of graphs, Discrete Math. Theor. Comput. Sci., 26 (2024), no. 3, Art. 16, 6 pp.
\bibitem[Petrusevski2018]{Petrusevski2018} M. Petru\v{s}evski, Odd $4$-edge-colorability of graphs, J. Graph Theory 87 (2018) 460--474.
\bibitem[Pyber1992]{Pyber1992} L. Pyber, Covering the edges of a graph by $\ldots$, in Sets, Graphs and Numbers (Budapest, 1991), Colloq. Math. Soc. J\'anos Bolyai 60, North-Holland, Amsterdam (1992) 583--610.
\bibitem[Scott1997]{Scott1997} A. D. Scott, On graph decompositions modulo $k$, Discrete Math. 175 (1997) 289--291.
\bibitem[Thomassen2014]{Thomassen2014} C. Thomassen, Graph factors modulo $k$, J. Combin. Theory Ser. B 106 (2014) 174--177.
\end{thebibliography}
\end{document}
